\section{Mathematical details: counting \texorpdfstring{$X_6$}{X6} and identifying boundary words}
\label{app:math_details}

\subsection{Counting admissible words}

Let $a_m$ be the number of length-$m$ binary words with no adjacent ones.
Classically,
\[
a_m=a_{m-1}+a_{m-2},\qquad a_1=2,\ a_2=3,
\]
since a legal word either starts with $0$ (followed by any legal word of length $m-1$) or starts with $10$ (followed by any legal word of length $m-2$).
Hence $a_6=21=F_8$.

\subsection{Boundary words at length 6}

A boundary word is a legal word $w=w_1\cdots w_6\in X_6$ with $w_1=w_6=1$.
The no-adjacent-ones constraint then forces $w_2=w_5=0$, and the remaining middle bits $(w_3,w_4)$ cannot be $(1,1)$.
Thus there are exactly three boundary words:
\[
100001,\qquad 100101,\qquad 101001.
\]

\subsection{Preimages under \texorpdfstring{$\mathrm{Fold}_6$}{Fold6}}
\label{app:fold6_preimages}

We record a complete preimage characterization for $\mathrm{Fold}_6$ restricted to $\{0,\dots,63\}$.
This makes the $m=6$ interface self-contained and yields closed-form proofs of Proposition~\ref{prop:boundary_preimages} and Theorem~\ref{thm:fold6_degeneracy}.

\begin{proposition}[Preimage characterization at window length $6$]
\label{prop:fold6_preimage_characterization}
Fix $w=(w_1,\dots,w_6)\in X_6$ and let $v:=V(w)\in\{0,\dots,20\}$.
Then every $N\in\{0,\dots,63\}$ satisfying $\mathrm{Fold}_6(N)=w$ has the form
\[
N=v+\delta,
\qquad
\delta\in\{0,21,34,55\}=\{0,F_8,F_9,F_{10}\},
\]
with the following admissibility constraints:
\begin{itemize}
\item if $w_6=1$ then necessarily $\delta\in\{0,34\}$;
\item if $w_6=0$ then $\delta\in\{0,21,34\}$ always, and $\delta=55$ is admissible if and only if $v\le 8$.
\end{itemize}
Consequently, the only possible preimage sizes are $2$, $3$, and $4$.
\end{proposition}

\begin{proof}
Let the Zeckendorf expansion of $N$ be
\[
N=\sum_{k\ge 1} c_k F_{k+1},
\qquad
c_k\in\{0,1\},
\qquad
c_k c_{k+1}=0.
\]
By definition $\mathrm{Fold}_6(N)=(c_1,\dots,c_6)=w$, hence
\[
N-V(w)=\sum_{k\ge 7} c_k F_{k+1}.
\]
Since $N\le 63$ and $F_{11}=89>63$, all digits beyond $k=9$ vanish and therefore
\[
\delta:=N-V(w)=c_7 F_8 + c_8 F_9 + c_9 F_{10}
 = 21c_7 + 34c_8 + 55c_9.
\]
The Zeckendorf constraint forbids adjacent ones, so $(c_7,c_8,c_9)$ contains at most one $1$ and thus $\delta\in\{0,21,34,55\}$.

If $w_6=c_6=1$ then adjacency forces $c_7=0$, hence $\delta\neq 21$.
Also $v\ge 13$ whenever $w_6=1$, so $v+55>63$ and $\delta\neq 55$.
Therefore $\delta\in\{0,34\}$.

If $w_6=c_6=0$ then $c_7$ may be $1$, so $\delta=21$ is admissible, and $\delta=34$ is always admissible (it only requires $c_8=1$ and $c_7=0$).
Finally, $\delta=55$ is admissible exactly when $v+55\le 63$, i.e.\ $v\le 8$.
\end{proof}

\begin{corollary}[Boundary preimages and the boundary index set]
The boundary words have Zeckendorf values $V(100001)=14$, $V(101001)=17$, $V(100101)=19$.
Since each boundary word has $w_6=1$, Proposition~\ref{prop:fold6_preimage_characterization} implies exactly two preimages differing by $34$:
\[
\mathrm{Fold}_6^{-1}(100001)=\{14,48\},\qquad
\mathrm{Fold}_6^{-1}(101001)=\{17,51\},\qquad
\mathrm{Fold}_6^{-1}(100101)=\{19,53\}.
\]
Hence
\[
\mathrm{Fold}_6^{-1}(X_6^{\mathrm{bdry}})=\{14,17,19,48,51,53\}.
\]
\end{corollary}

\begin{corollary}[Degeneracy histogram for $\mathrm{Fold}_6$ on $\{0,\dots,63\}$]
The degeneracy counts in Theorem~\ref{thm:fold6_degeneracy} follow from Proposition~\ref{prop:fold6_preimage_characterization}:
\[
\begin{aligned}
\#\{w\in X_6:\ |\mathrm{Fold}_6^{-1}(w)|=2\} &= 8,\\
\#\{w\in X_6:\ |\mathrm{Fold}_6^{-1}(w)|=3\} &= 4,\\
\#\{w\in X_6:\ |\mathrm{Fold}_6^{-1}(w)|=4\} &= 9.
\end{aligned}
\]
\end{corollary}

\begin{proof}
If $w_6=1$ then Proposition~\ref{prop:fold6_preimage_characterization} gives $|\mathrm{Fold}_6^{-1}(w)|=2$.
The constraint $w_6=1$ forces $w_5=0$, and the remaining prefix $(w_1,\dots,w_4)$ is any admissible length-$4$ word.
Thus there are $a_4=F_6=8$ such words.

If $w_6=0$ and $v>8$, then $|\mathrm{Fold}_6^{-1}(w)|=3$ (the $\delta=55$ option is excluded by $v+55>63$).
For admissible $w$ with $w_6=0$, one has $v>8$ if and only if $w_5=1$ and at least one of $(w_1,w_2,w_3)$ is $1$ (since $w_4$ must be $0$ by admissibility).
The length-$3$ admissible prefixes are exactly $\{000,100,010,001,101\}$, so there are $4$ nonzero choices and hence $4$ such words.

All remaining admissible words have $w_6=0$ and $v\le 8$, hence have preimage size $4$.
Since $|X_6|=21$, the remaining count is $21-8-4=9$.
\end{proof}


